\chapter{事件与图件索引}
\label{app:index}

\section{索引使用}
\label{app:index:guide}

本附录只收录已在正文设置跳转目标的事件和已置入正文的图件；事件账本不是索引的替代品。事件名称一律使用中文可点击链接，内部 ID 只用于锚点和维护，不向读者显示。材料的性质、适用范围和主要限制见\hyperref[app:sources]{“来源与史料说明”}。

\section{事件索引}
\label{app:index:events}

事件按西汉与新莽、东汉、三国、西晋、东晋与十六国、南北朝分组编列。新增条目必须先在正文首次展开处设置 \verb|\eventanchor{内部ID}{中文名称}|，再在相应分组以 \verb|\eventindexentry{内部ID}{中文名称}| 登记；尚无正文锚点的账本事件不在此建立失效链接。

\begin{description}
  \item[西汉与新莽] 正文事件锚点建立后在此登记。
  \item[东汉] 正文事件锚点建立后在此登记。
  \item[三国] 正文事件锚点建立后在此登记。
  \item[西晋] 正文事件锚点建立后在此登记。
  \item[东晋与十六国] 正文事件锚点建立后在此登记。
  \item[南北朝] 正文事件锚点建立后在此登记。
\end{description}

\section{图件索引}
\label{app:index:figures}

图件按时代分组。新增图件须先在图件环境或图注处设置稳定的 \verb|\label{fig:...}|，再以 \verb|\figureindexentry{fig:...}{图名}| 登记；没有标签的图件不在此建立链接。

\begin{description}
  \item[西汉与新莽] 已标注图件在此登记。
  \item[东汉] 已标注图件在此登记。
  \item[三国] 已标注图件在此登记。
  \item[西晋] 已标注图件在此登记。
  \item[东晋与十六国] 已标注图件在此登记。
  \item[南北朝] 已标注图件在此登记。
\end{description}

\section{来源入口}
\label{app:index:sources}

按材料类型、时段和使用边界查找来源，请参看\hyperref[app:sources]{“来源与史料说明”}；其中的\hyperref[app:sources:narrative]{“基础叙事与编年材料”}、\hyperref[app:sources:documents]{“文书、宗教与物质材料”}和\hyperref[app:sources:periods]{“按时段进入材料”}可与本索引配合使用。
